% =============================================================================
% Capítulo 1 — Introdução
% Começar cada capítulo numa nova página e incluir uma breve descrição
% sobre qual será o assunto do capítulo.
% =============================================================================

\chapter{Introdução}
\label{cap:introducao}

Este capítulo apresenta o contexto e a motivação do trabalho, a definição do problema e os objetivos propostos, a solução implementada e a organização do presente documento.

% ---- 1.1 Contexto e Motivação ----
\section{Contexto e Motivação}
\label{sec:contexto_motivacao}

% Responder às seguintes questões:
% - Em que contexto foi desenvolvido o projeto?
% - Qual é a área ou áreas científicas do projeto?
% - Quais os agentes envolvidos (empresa, laboratório de investigação, etc.)?
% - O que motivou a realização do projeto?
Descrever o contexto geral do trabalho, a área em que se insere, os agentes envolvidos e a sua relevância.

% ---- 1.2 Definição do Problema e Objetivos ----
\section{Definição do Problema e Objetivos}
\label{sec:problema_objetivos}

% Responder às seguintes questões:
% - Qual foi o problema que se pretendeu resolver?
% - Quais foram os objetivos do projeto?
% Tente identificar objetivos mensuráveis.
Identificar e descrever o problema que motivou a realização deste trabalho.

Apresentar uma lista formal dos objetivos do projeto. Tipicamente, os objetivos correspondem às funcionalidades que se pretende que o sistema disponibilize ou ao conhecimento que se pretende adquirir:

\begin{itemize}
    \item Objetivo 1
    \item Objetivo 2
    \item Objetivo 3
\end{itemize}

% ---- 1.3 Solução Proposta ----
\section{Solução Proposta}
\label{sec:solucao}

% Responder à seguinte questão:
% - Qual foi a solução implementada?
% Apresentar uma visão geral da solução (o sistema desenvolvido de uma
% perspetiva de alto nível), mencionando as principais tecnologias usadas.
Descrever brevemente a solução proposta para o problema identificado.

% ---- 1.4 Organização do Relatório ----
\section{Organização do Relatório}
\label{sec:organizacao}

% Responder à seguinte questão:
% - Como está o relatório organizado?
% Escrever um parágrafo para cada capítulo descrevendo o seu assunto.
O presente documento encontra-se organizado em cinco capítulos:

\begin{itemize}
    \item \textbf{Capítulo~\ref{cap:introducao} --- Introdução:} Apresenta o contexto, a motivação, o problema, os objetivos e a solução proposta.
    \item \textbf{Capítulo~\ref{cap:trabalho_relacionado} --- Trabalho Relacionado:} Analisa os métodos, tecnologias e soluções existentes relevantes para o projeto.
    \item \textbf{Capítulo~\ref{cap:desenvolvimento} --- Desenvolvimento:} Descreve a análise de requisitos, a arquitetura do sistema e a implementação.
    \item \textbf{Capítulo~\ref{cap:testes} --- Testes e Resultados:} Apresenta os testes realizados e os resultados obtidos.
    \item \textbf{Capítulo~\ref{cap:conclusoes} --- Conclusões:} Resume as conclusões e sugere trabalho futuro.
\end{itemize}
